\documentclass[11pt,a4paper]{article}

\usepackage{ctex}
\usepackage[T1]{fontenc}
\usepackage[top=2.5cm, bottom=2.5cm, left=2.5cm, right=2.5cm]{geometry}
\usepackage{booktabs}
\usepackage{array}
\usepackage{enumitem}
\usepackage{tabularx}
\usepackage{amssymb}
\usepackage{xcolor}
\definecolor{primary}{HTML}{1a1a2e}
\definecolor{accent}{HTML}{3b82f6}
\definecolor{success}{HTML}{22c55e}
\definecolor{danger}{HTML}{ef4444}
\definecolor{warning}{HTML}{f59e0b}
\definecolor{graytext}{HTML}{6b7280}
\definecolor{progress}{HTML}{8b5cf6}
\usepackage[hidelinks]{hyperref}
\hypersetup{colorlinks=true, linkcolor=primary, urlcolor=accent}
\usepackage{titlesec}
\titleformat{\section}{\Large\bfseries\color{primary}}{}{0em}{}[\titlerule]
\titleformat{\subsection}{\large\bfseries\color{primary}}{}{0em}{}
\titleformat{\subsubsection}{\normalsize\bfseries\color{primary}}{}{0em}{}
\usepackage{fancyhdr}
\pagestyle{fancy}
\fancyhf{}
\fancyhead[L]{\small\color{graytext} 企业级 AI Agent 应用商店 -- 项目方法论}
\fancyhead[R]{\small\color{graytext} GUIDE-002}
\fancyfoot[C]{\small\color{graytext} \thepage}
\renewcommand{\headrulewidth}{0.4pt}

\title{\textbf{GUIDE-002：项目方法论总纲}}
\author{万能产品小助手（PM AI）\\魏子政 审阅}
\date{2026-05-06 / 版本 v1.0}

\begin{document}

\maketitle

\begin{center}
\begin{tabular}{>{\bfseries}p{3cm} p{10cm}}
\toprule
提出日期 & 2026-05-06 \\
提出者 & 魏子政 \\
优先级 & \textcolor{danger}{P0} \\
状态 & \textcolor{success}{已完成} \\
文档类型 & 项目方法论（永久维护） \\
\bottomrule
\end{tabular}
\end{center}

\vspace{1em}

\section{文档定位}

本文档是项目的方法论总纲，回答以下核心问题：

\begin{enumerate}[nosep]
  \item 两种智能体（OpenClaw + Cursor）的选型和搭配细节是什么？
  \item 属于 OpenClaw 的 Skill 和属于 Cursor 的 Skill 分别怎么加载？有何区别？
  \item 每个技术组件为什么选它而不选竞品？
  \item 项目需求架构文档（tex-pdf）如何维护和撰写？
  \item 后续需求新增如何处理？
  \item 文档版本如何迭代？
  \item 自闭环机制如何运行？
\end{enumerate}

本文档整合自 AGENT-ARCHITECTURE.md、tech-selection.md、DOC-LIFECYCLE.md、TEAM-MANUAL.md 等多份独立文档，作为唯一的方法论入口。

---

\section{智能体搭配架构}

\subsection{为什么需要两种智能体}

企业级应用需要两类能力：\textbf{产品管理}（需求理解、任务拆解、质量审计、进度追踪）和\textbf{代码开发}（编写前后端代码、运行命令、构建项目）。

\begin{table}[h]
\centering
\begin{tabularx}{\linewidth}{lX}
\toprule
\textbf{能力} & \textbf{说明} \\
\midrule
产品管理 & 需求理解、任务拆解、质量审计、进度追踪、文档维护、飞书通信 \\
代码开发 & 编写前后端代码、运行终端命令、理解项目上下文、构建部署 \\
\bottomrule
\end{tabularx}
\end{table}

\textbf{核心分工}：魏子政只跟 OpenClaw 对话 $\to$ OpenClaw 拆解需求并调度 Cursor $\to$ Cursor 写代码 $\to$ OpenClaw 审计质量 $\to$ 向魏子政汇报。

\subsection{两种智能体的本质区别}

\begin{table}[h]
\centering
\begin{tabularx}{\linewidth}{llX}
\toprule
\textbf{维度} & \textbf{OpenClaw (PM)} & \textbf{Cursor (开发)} \\
\midrule
定位 & 常驻 AI 助手 + PM & 代码编辑器内置 AI \\
记忆 & 文件级记忆（memory/、MEMORY.md） & 无跨会话记忆 \\
通信 & 飞书通道，可主动推送 & 无外部通信能力 \\
 Skill 来源 & \textasciitilde/.openclaw/skills/ + 工作区 skills/ & CLAUDE.md + .cursor/rules/*.mdc \\
Skill 加载 & 按需 read SKILL.md（AI 自主判断） & 自动扫描加载（alwaysApply / globs） \\
Skill 格式 & 普通 Markdown & MDC（YAML front matter + Markdown） \\
触发方式 & 魏子政发消息到飞书 & OpenClaw 通过 CLI 调度 \\
调用方式 & N/A（自己是调度者） & cursor agent --print --trust --workspace <路径> \\
能力边界 & 不直接改代码（代码走 Cursor） & 不飞书通信、无持久记忆 \\
\bottomrule
\end{tabularx}
\end{table}

\subsection{Skill 加载机制详解}

\subsubsection{OpenClaw 的 Skill 加载}

\begin{enumerate}[nosep]
  \item OpenClaw 启动会话时，系统提示中包含 \texttt{<available\_skills>} 列表（每个 skill 的名称和 SKILL.md 路径）
  \item 收到用户消息后，OpenClaw \textbf{根据消息内容自动判断}应该加载哪个 skill
  \item 用 \texttt{read} 工具读取对应的 SKILL.md 文件内容，作为后续行动的指令
\end{enumerate}

存放位置：
\begin{itemize}[nosep]
  \item \textasciitilde/.openclaw/skills/ --- 全局 skills（所有项目共享）
  \item 工作区 skills/ --- 项目专属 skills
\end{itemize}

\subsubsection{Cursor 的规则加载}

\begin{enumerate}[nosep]
  \item \texttt{CLAUDE.md} --- 项目根目录，Cursor \textbf{每次对话自动加载}，作为总入口
  \item \texttt{.cursor/rules/*.mdc} --- MDC 格式规则文件，支持两种触发：
  \begin{itemize}[nosep]
    \item \texttt{alwaysApply: true} --- 每次对话都加载
    \item \texttt{globs: ["backend/**"]} --- 只在编辑匹配目录下的文件时加载
  \end{itemize}
\end{enumerate}

MDC 文件格式：
\begin{verbatim}
---
description: 规则描述
globs: ["backend/**"]     ← 可选
alwaysApply: true          ← 可选
---
# 规则标题
规则内容...
\end{verbatim}

\subsubsection{本项目中的 Skill/规则清单}

\begin{table}[h]
\centering
\begin{tabularx}{\linewidth}{llp{6cm}}
\toprule
\textbf{文件} & \textbf{给谁} & \textbf{用途} \\
\midrule
\texttt{skills/pm/SKILL.md} & OpenClaw & 需求拆解、任务分配、阶段审计、查证 \\
\texttt{skills/ui-design/SKILL.md} & OpenClaw & 设计方案输出、视觉规范把控 \\
\texttt{skills/tester/SKILL.md} & OpenClaw & 阶段门控验证、功能测试 \\
\texttt{CLAUDE.md} & Cursor & 每次必加载的项目总入口 \\
\texttt{.cursor/rules/project-overview.mdc} & Cursor & 技术栈锁定、禁止事项 \\
\texttt{.cursor/rules/backend-rules.mdc} & Cursor & 后端编码规范 \\
\texttt{.cursor/rules/frontend-rules.mdc} & Cursor & 前端编码规范 \\
\texttt{.cursor/rules/ui-design.mdc} & Cursor & UI 设计规范 \\
\texttt{.cursor/rules/api-conventions.mdc} & Cursor & API 接口规范 \\
\texttt{.cursor/rules/testing.mdc} & Cursor & 测试自检要求 \\
\texttt{.cursor/rules/database.mdc} & Cursor & 数据库编码约定 \\
\bottomrule
\end{tabularx}
\end{table}

\textbf{关键理解}：\texttt{skills/} 目录下的 SKILL.md 是给 \textbf{OpenClaw} 看的（让 PM 知道规范长什么样），\texttt{.cursor/rules/} 是给 \textbf{Cursor} 看的（让开发时自动遵守规范）。两者内容有关联但不相同。

\subsection{调用链路}

一次完整任务执行的完整链路：

\begin{enumerate}[nosep]
  \item \textbf{魏子政}在飞书发消息提出需求
  \item \textbf{OpenClaw} 加载 pm/SKILL.md，读取 memory 文件恢复上下文，读取项目文档
  \item OpenClaw 执行 PM 职责：记录任务 $\to$ 拆解需求 $\to$ 更新 PROGRESS.md
  \item OpenClaw 构造 CLI 命令调度 Cursor：
  \begin{verbatim}
  cursor agent --print --trust \
    --workspace <项目路径> "任务描述"
  \end{verbatim}
  \item \textbf{Cursor} 自动加载 CLAUDE.md + rules/，读取相关文件，编写代码，返回结果
  \item \textbf{OpenClaw} 加载 tester/SKILL.md，审计 Cursor 产出（查证三问），如有问题退回修复
  \item OpenClaw 向魏子政汇报，更新所有受影响文档
\end{enumerate}

---

\section{技术选型与理由}

\subsection{锁定技术栈总览}

\begin{table}[h]
\centering
\begin{tabularx}{\linewidth}{lX}
\toprule
\textbf{层} & \textbf{技术} \\
\midrule
前端 & Vue 3 + Vite + Element Plus + Pinia + TypeScript \\
后端 & FastAPI + SQLAlchemy (async) + Pydantic \\
数据库 & PostgreSQL 16（JSONB 存扫描结果） \\
对象存储 & MinIO（S3 兼容） \\
搜索引擎 & OpenSearch 2.x（全文 + 向量） \\
认证授权 & JWT (python-jose + passlib) + RBAC \\
容器化 & Docker + Docker Compose \\
数据库迁移 & Alembic \\
\bottomrule
\end{tabularx}
\end{table}

\textbf{硬约束}：前端不用 Next.js/React，后端不用 Spring/Java。

\subsection{逐项选型理由}

\subsubsection{Vue 3 + Element Plus（前端）}

\begin{itemize}[nosep]
  \item 学习曲线低，模板语法直观，实习生上手快
  \item Element Plus 企业级组件库开箱即用（表格/表单/弹窗/分页），减少 50\%+ UI 工作量
  \item 国内团队主流，中文文档完善，社区活跃
  \item TypeScript 一等支持，\texttt{<script setup>} + TS 组合式 API
  \item Vite HMR 毫秒级，开发体验好
\end{itemize}

\textbf{不用 Next.js/React}：SSR/SSG 对内部管理系统价值不大，多出的复杂度是负担；React 额外状态管理方案（Redux/Zustand），Vue 的 Pinia 更简洁。

\subsubsection{FastAPI（后端）}

\begin{itemize}[nosep]
  \item async/await 原生支持，天然适合 IO 密集型（上传/搜索/扫描）
  \item Pydantic 自动生成 Swagger/OpenAPI 文档，前后端对接零成本
  \item Python 生态与安全扫描工具（Semgrep/ClamAV）、LLM API 天然互通
  \item 代码量少，原型快，5-10 人团队产出高
\end{itemize}

\textbf{不用 Spring Boot}：Java 启动慢、代码量大，与 Python 安全扫描工具集成不便。不用 NestJS：团队主语言 Python，减少跨语言成本。

\subsubsection{PostgreSQL 16（数据库）}

\begin{itemize}[nosep]
  \item ACID 保证，适合元数据、权限、审计等强一致性数据
  \item JSONB 类型存储扫描结果（嵌套 JSON），支持查询和索引
  \item 扩展性强：pgvector、全文检索、窗口函数、CTE
  \item 运维成熟，免费开源
\end{itemize}

\textbf{不用 MongoDB}：减少组件数，JSONB 足够。\textbf{不用 MySQL}：JSONB 支持不如 PG。

\subsubsection{OpenSearch 2.x（搜索）}

\begin{itemize}[nosep]
  \item 全文 + 向量一体化，kNN 语义搜索
  \item DSL 支持 per-document 权限查询
  \item IK Analysis 中文分词
  \item Apache 2.0 许可（无 SSPL 限制）
\end{itemize}

\textbf{不用 Elasticsearch}：许可证变更（SSPL），商业使用有合规风险。

\subsubsection{MinIO（对象存储）}

\begin{itemize}[nosep]
  \item S3 API 兼容，随时可切换到 AWS S3/阿里 OSS
  \item 私有化部署，数据不出域
  \item 预签名 URL，前端直接下载不经后端转发
  \item Docker 一容器搞定，自带 Web 管理控制台
\end{itemize}

\subsubsection{JWT + RBAC（认证授权）}

\begin{itemize}[nosep]
  \item 轻量：Keycloak 需要 JVM + 独立服务，JWT 只需两个 Python 库
  \item 够用：初期 4 个角色（admin/reviewer/developer/viewer）
  \item 无状态：Token 自包含，不需要每次查库验证
\end{itemize}

\textbf{扩展路径}：需要 LDAP/AD/SSO 时再引入 Keycloak，JWT 签发逻辑改对接即可。

\subsubsection{Docker Compose（容器化）}

\begin{itemize}[nosep]
  \item MVP 不需要 K8s，单机部署足够
  \item \texttt{docker compose up} 一键启动所有依赖，新人 5 分钟搭环境
  \item 5-10 人团队不需要专门 K8s 运维人员
\end{itemize}

\textbf{扩展路径}：需要多实例/灰度发布时迁移到 K8s（Dockerfile 不改，加 Helm Chart）。

\subsection{技术栈锁定规则}

\begin{enumerate}[nosep]
  \item 技术栈一经锁定，OpenClaw 和 Cursor 均不得自行变更
  \item 如需变更：魏子政确认 $\to$ OpenClaw 评估影响 $\to$ 更新文档 $\to$ 更新 Cursor rules
  \item 新增依赖（npm/pip 包）允许，但不得引入新框架/数据库/中间件
  \item 所有变更记录在 tech-selection.md 变更记录中
\end{enumerate}

---

\section{文档体系与维护}

\subsection{文档目录结构}

\begin{verbatim}
enterprise-app-store/
├── 根目录（方法论与规范）
│   ├── AGENT-ARCHITECTURE.md   ← 智能体搭配（详细版）
│   ├── tech-selection.md       ← 技术选型（详细版）
│   ├── TEAM-MANUAL.md          ← 团队协作指南
│   ├── DOC-LIFECYCLE.md        ← 文档生命周期（详细版）
│   ├── CHANGELOG.md            ← 统一变更日志
│   ├── PROGRESS.md             ← 进度追踪
│   ├── CLAUDE.md               ← Cursor 总入口
│   └── BUGS.md                 ← Bug 记录
│
├── core/（核心产品文档链）
│   ├── MRD.tex/pdf             ← 市场需求文档
│   ├── PRD.tex/pdf             ← 产品需求文档
│   ├── RFC.tex/pdf             ← 设计方案
│   └── TRD.tex/pdf             ← 技术设计文档
│
├── demands/（需求文档）
│   ├── MASTER.tex/pdf          ← 活跃需求 Backlog（持续更新）
│   ├── DEMAND-xxx.tex/pdf      ← 冻结的需求快照
│   └── GUIDE-xxx.tex/pdf       ← 项目指南
│
├── reference/（外部研究）
│   ├── REF-xxx.tex/pdf         ← 外部架构调研
│   └── DOC-NAMING-CONVENTION   ← 文档命名规范索引
│
├── .cursor/rules/*.mdc         ← Cursor 编码规则
└── skills/xxx/SKILL.md         ← OpenClaw 角色 Skill
\end{verbatim}

\subsection{文档分类与维护责任}

\begin{table}[h]
\centering
\begin{tabularx}{\linewidth}{llX}
\toprule
\textbf{分类} & \textbf{维护者} & \textbf{更新时机} \\
\midrule
根目录方法论 & OpenClaw & 需求变更、技术决策、流程调整 \\
core/ 核心文档 & OpenClaw & 功能模块变更、架构调整 \\
demands/ 需求 & OpenClaw & 新需求提出、需求拆解 \\
reference/ 研究 & OpenClaw & 外部调研完成 \\
Cursor rules & OpenClaw & 编码规范变更 \\
代码 & Cursor & 功能开发、Bug 修复 \\
\bottomrule
\end{tabularx}
\end{table}

\subsection{tex-pdf 生成流程}

\begin{enumerate}[nosep]
  \item 使用 Tectonic LaTeX 引擎：\texttt{\textasciitilde/.local/tectonic/tectonic input.tex}
  \item 必须先 \texttt{cd} 到 tex 文件所在目录再执行（相对路径依赖）
  \item tex 是源文件（Git 管理），pdf 是交付物（也提交 Git 方便查看）
  \item 已知限制：不用 fontawesome5 命令、不用 Unicode 箭头（用 \texttt{\textbackslash to}）、\texttt{\textbackslash square} 需数学模式
\end{enumerate}

\subsection{文档版本管理}

采用语义化版本 \texttt{vX.Y.Z}：

\begin{table}[h]
\centering
\begin{tabularx}{\linewidth}{lX}
\toprule
\textbf{数字} & \textbf{含义} \\
\midrule
X（主版本） & 架构重大变更、技术栈变更 \\
Y（次版本） & 新增模块/功能、重要规范新增 \\
Z（补丁版本） & 文档修正、措辞优化、小补丁 \\
\bottomrule
\end{tabularx}
\end{table}

每个重要文档头部包含版本号、日期、更新人、变更摘要。

---

\section{需求新增处理流程}

\subsection{需求来源}

\begin{table}[h]
\centering
\begin{tabularx}{\linewidth}{lX}
\toprule
\textbf{来源} & \textbf{处理方式} \\
\midrule
魏子政直接提 & 记录 $\to$ 评估 $\to$ 拆解 $\to$ 排期 \\
阶段审计发现 & 记录 $\to$ 评估 $\to$ 排入下阶段 \\
Bug 修复衍生 & 随 Bug 修复一起处理 \\
\bottomrule
\end{tabularx}
\end{table}

\subsection{需求处理五步流程}

\textbf{第 1 步：需求记录}

写入 \texttt{memory/YYYY-MM-DD.md}，包含需求描述、提出者、优先级。复杂需求追问魏子政确认细节。

\textbf{第 2 步：影响评估（OpenClaw 执行）}

\begin{itemize}[nosep]
  \item 技术影响：需要改哪些模块？新增哪些表/接口/页面？
  \item 文档影响：哪些文档需要更新？
  \item 进度影响：对当前里程碑有无影响？
\end{itemize}

\textbf{第 3 步：方案设计（复杂需求）}

\begin{itemize}[nosep]
  \item 新增数据库表？$\to$ 更新 SCHEMA
  \item 新增 API？$\to$ 补充 developer SKILL
  \item 涉及技术栈变更？$\to$ \textbf{必须魏子政审批}
  \item 方案输出后给魏子政确认
\end{itemize}

\textbf{第 4 步：任务拆解与分配}

\begin{itemize}[nosep]
  \item 拆解为可独立完成和验证的开发任务
  \item 写入 PROGRESS.md
  \item 分配给 Cursor 执行
\end{itemize}

\textbf{第 5 步：开发 $\to$ 测试 $\to$ 审计 $\to$ 交付}

\begin{itemize}[nosep]
  \item Cursor 执行开发
  \item OpenClaw 独立验收（查证三问）
  \item 更新所有受影响文档
  \item 向魏子政汇报
\end{itemize}

\subsection{需求优先级}

\begin{table}[h]
\centering
\begin{tabularx}{\linewidth}{llX}
\toprule
\textbf{优先级} & \textbf{含义} & \textbf{处理时效} \\
\midrule
P0 & 阻塞当前阶段推进 & 立即处理 \\
P1 & 核心功能缺陷/缺失 & 当天处理 \\
P2 & 重要但不紧急 & 排入当前/下阶段 \\
P3 & 锦上添花 & 有空再做 \\
\bottomrule
\end{tabularx}
\end{table}

\subsection{需求变更控制}

\textbf{小变更}（不影响已有模块）：OpenClaw 直接处理，事后汇报。

\textbf{中变更}（影响已有模块但不涉及架构）：OpenClaw 评估 $\to$ 简要说明 $\to$ 执行。

\textbf{大变更}（涉及架构/技术栈/数据库结构）：OpenClaw 评估 $\to$ 详细方案 $\to$ \textbf{魏子政审批} $\to$ 执行。

\subsection{需求文档维护：MASTER + DEMAND 双层体系}

\begin{itemize}[nosep]
  \item \textbf{MASTER.tex}（活文档）：持续更新的需求 Backlog，所有需求的最新状态
  \item \textbf{DEMAND-xxx.tex}（冻结快照）：每个需求冻结后的完整版本，不再修改
  \item 新需求提出时：先写入 MASTER（分配编号 D-xxx），开发完成后冻结为 DEMAND-xxx
  \item MASTER 是唯一权威来源，DEMAND 是历史快照
\end{itemize}

---

\section{自闭环机制}

\subsection{什么是自闭环}

自闭环 = 从需求提出到交付完成，\textbf{所有环节在同一体系内完成}，不需要外部干预。

\begin{verbatim}
需求输入（魏子政）
    ↓
OpenClaw 接收 → 记录 → 评估 → 设计
    ↓
Cursor 开发 → 代码产出
    ↓
OpenClaw 独立验收 → 查证
    ↓
文档同步更新 → Git commit
    ↓
交付输出（魏子政）
    ↓
反馈 → 回到 OpenClaw → 循环
\end{verbatim}

\subsection{五个闭环}

\textbf{闭环 1：需求闭环}
\begin{itemize}[nosep]
  \item 输入：魏子政提需求
  \item 过程：记录 $\to$ 评估 $\to$ 设计 $\to$ 拆解
  \item 闭合标志：PROGRESS.md 出现新任务，MASTER.tex 更新
\end{itemize}

\textbf{闭环 2：开发闭环}
\begin{itemize}[nosep]
  \item 输入：Cursor 收到任务
  \item 过程：读 rules/ $\to$ 写代码 $\to$ 自测
  \item 闭合标志：\texttt{npm run build} 通过 + 后端启动正常
\end{itemize}

\textbf{闭环 3：验收闭环}
\begin{itemize}[nosep]
  \item 输入：Cursor 产出代码
  \item 过程：OpenClaw 独立验收（查证三问 + 功能验证）
  \item 闭合标志：验收清单全部通过
\end{itemize}

\textbf{闭环 4：文档闭环}
\begin{itemize}[nosep]
  \item 输入：代码变更 / 需求变更
  \item 过程：更新受影响文档 $\to$ 生成新 PDF $\to$ 记录 CHANGELOG
  \item 闭合标志：CHANGELOG.md 有新记录，文档版本号已更新
\end{itemize}

\textbf{闭环 5：审计闭环}
\begin{itemize}[nosep]
  \item 输入：验收通过的开发任务
  \item 过程：查证三问（目标达成？细节硬伤？换位质疑？）
  \item 闭合标志：查证三问全部通过，Git commit 完成
\end{itemize}

\subsection{闭环监控规则}

\textbf{如果任何闭环未闭合，任务不能标记为「已完成」。}

\begin{table}[h]
\centering
\begin{tabularx}{\linewidth}{lX}
\toprule
\textbf{检查项} & \textbf{闭合标准} \\
\midrule
需求闭环 & PROGRESS.md 有任务记录 + MASTER.tex 已更新 \\
开发闭环 & Cursor 返回成功 + build 通过 \\
验收闭环 & 验收清单全部通过 \\
文档闭环 & 受影响文档已更新 + CHANGELOG.md 有记录 \\
审计闭环 & 查证三问通过 + Git commit 完成 \\
\bottomrule
\end{tabularx}
\end{table}

\subsection{查证三问（每次验收必执行）}

\begin{enumerate}[nosep]
  \item \textbf{目标真的达成了吗？} --- 产出物是否完整覆盖原始需求？
  \item \textbf{细节有硬伤吗？} --- 技术细节、数据、路径、命令是否可验证为正确？
  \item \textbf{如果我是接收者，我会问什么？} --- 换位思考，发散性质疑
\end{enumerate}

---

\section{文档版本迭代}

\subsection{更新触发条件}

\begin{table}[h]
\centering
\begin{tabularx}{\linewidth}{lX}
\toprule
\textbf{触发条件} & \textbf{需要更新的文档} \\
\midrule
技术栈变更 & tech-selection.md $\to$ 重新生成相关 PDF \\
新增功能模块 & core/ 文档 + demands/ DEMAND-xxx \\
数据库结构变更 & TRD.tex + SCHEMA \\
需求新增 & MASTER.tex + DEMAND-xxx.tex \\
规范变更 & rules/*.mdc + SKILL.md \\
阶段里程碑完成 & core/ 文档版本号更新 \\
\bottomrule
\end{tabularx}
\end{table}

\subsection{文档更新流程}

\begin{enumerate}[nosep]
  \item 触发条件出现
  \item OpenClaw 判断哪些文档需要更新
  \item 更新 Markdown/tex 源文件
  \item 如果是 tex 文档 $\to$ Tectonic 生成新 PDF
  \item 更新文档版本号（头部 vX.Y.Z）
  \item 记录到 CHANGELOG.md
  \item 更新 DOC-NAMING-CONVENTION 索引
  \item 更新 PROGRESS.md
  \item 向魏子政汇报（变更摘要 + 新文档位置）
\end{enumerate}

\subsection{Git 提交规范}

\begin{verbatim}
feat: 新增应用评分功能
fix: 修复登录页 token 过期未跳转的问题
docs: 更新技术选型文档 v2.0
refactor: 重构上传模块为异步
\end{verbatim}

文档变更和对应代码变更放在同一个 commit 或紧邻的 commit，不允许文档与代码不同步。

---

\section{分工矩阵}

\begin{table}[h]
\centering
\begin{tabularx}{\linewidth}{llX}
\toprule
\textbf{工作} & \textbf{谁做} & \textbf{怎么做} \\
\midrule
接收需求 & OpenClaw & 飞书对话 \\
需求拆解 & OpenClaw & 读取 pm/SKILL.md \\
架构设计 & OpenClaw & 查证技术文档 \\
数据库设计 & OpenClaw 设计 $\to$ Cursor 实现 & SCHEMA + database.mdc \\
UI 设计 & Cursor & 根据 ui-design.mdc 规范自主设计实现 \\
写后端代码 & Cursor & backend-rules.mdc + api-conventions.mdc \\
写前端代码 & Cursor & frontend-rules.mdc + ui-design.mdc \\
环境搭建 & Cursor & Docker Compose \\
自验收 & Cursor & testing.mdc \\
独立验收 & OpenClaw & 查证三问 \\
修 Bug & Cursor + OpenClaw & BUGS.md 记录 \\
写文档 & OpenClaw & tex-pdf + Markdown \\
生成 PDF & OpenClaw & Tectonic \\
进度管理 & OpenClaw & PROGRESS.md \\
汇报沟通 & OpenClaw & 飞书 + Edge TTS 语音 \\
需求变更评估 & OpenClaw & 分析影响 + 更新文档 \\
Git commit & OpenClaw & 统一提交 \\
\bottomrule
\end{tabularx}
\end{table}

---

\section{更新记录}

\begin{tabular}{lll}
\toprule
版本 & 日期 & 变更内容 \\
\midrule
v1.0 & 2026-05-06 & 初始版本，整合方法论总纲 \\
\bottomrule
\end{tabular}

\end{document}
